% Overleaf-ready section fragment.
% Recommended packages in the parent document:
% \usepackage{hyperref}
% \usepackage{booktabs}
% Insert with:
% % Overleaf-ready section fragment.
% Recommended packages in the parent document:
% \usepackage{hyperref}
% \usepackage{booktabs}
% Insert with:
% % Overleaf-ready section fragment.
% Recommended packages in the parent document:
% \usepackage{hyperref}
% \usepackage{booktabs}
% Insert with:
% % Overleaf-ready section fragment.
% Recommended packages in the parent document:
% \usepackage{hyperref}
% \usepackage{booktabs}
% Insert with:
% \input{TESSARIS_DEPARTMENT_PILOTS_AND_MARKETING_CREATIVE_STACK_2026-08-10}

\section{Unified Department Pilots and Marketing Creative Production Stack}
\label{sec:tessaris-department-pilots-marketing-creative}

\subsection{Scope and design objective}

This development phase consolidated the Tessaris department workspaces and
advanced Marketing from an early social-post prototype into a governed creative
production and experimentation system. The work covered five connected layers:

\begin{enumerate}
    \item a consistent department page structure for Marketing, Sales, Finance,
    Operations, Support and Human Resources;
    \item department-scoped AION/Pilot conversations and Boardroom assignment
    packages;
    \item consolidation of the former standalone Marketing Stream into the
    Marketing Pilot workspace;
    \item a provider-neutral Creative Intelligence Lab for images, video, clips,
    platform renditions and controlled experiments; and
    \item exact-approved paid video rendering plus locally owned timeline and
    clip export.
\end{enumerate}

The central architectural principle is that Tessaris owns the business context,
creative direction, governance, asset provenance and learning loop. External
generative-AI providers remain replaceable production engines. This allows the
product to perform as much work as possible locally while using specialist or
premium services only where model-scale generation is genuinely required.

\subsection{Unified department workspace shell}

The earlier department pages had diverged. Some routes rendered a full-screen
spatial Boardroom, some retained an obsolete black Pilot terminal, and some hid
the newer operational workspace. Support and Human Resources were especially
affected because legacy containers could overlap or replace the newer interfaces.

Every principal department now follows the same hierarchy:

\begin{enumerate}
    \item \textbf{Compact department director.} A small department-specific AION
    character identifies the function and its advisory boundary. It replaces the
    oversized spatial room on operational pages.
    \item \textbf{Department-scoped conversation.} The \emph{Talk to Pilot}
    control opens a conversation whose context is limited to the active
    department. For example, the Support Pilot discusses customer cases,
    policies and escalations rather than acting as a generic assistant.
    \item \textbf{Boardroom assignment package.} The package is retained because
    it contains the Boardroom meeting instructions, goals, discovery requirements,
    linked goal sheet and expected outputs. The \emph{Open Package} control is
    available consistently.
    \item \textbf{Operational workspace.} The real department system appears
    beneath the package: Sales Revenue, Finance, Support Case Centre, HR people
    and permissions, Marketing, or the Operations foundation.
\end{enumerate}

Legacy black Pilot containers and oversized spatial canvases are explicitly
excluded from these routes when they would override the operational workspace.
The installed department outcomes are:

\begin{table}[htbp]
\centering
\small
\begin{tabular}{p{0.17\linewidth} p{0.70\linewidth}}
\toprule
\textbf{Department} & \textbf{Mounted operational workspace} \\
\midrule
Marketing & Unified automation, creative production, approvals, calendar and history \\
Sales & Canonical CRM, governed communications and Sales Revenue workspace \\
Finance & Finance operating workspace, evidence, bookkeeping and reporting controls \\
Operations & Unified department shell and foundation for the Operations build \\
Support & Canonical Support Case Centre, policy authority and escalation workflow \\
Human Resources & Organisation, people, roles, permissions and authority workspace \\
\bottomrule
\end{tabular}
\caption{Standardised department workspace structure.}
\end{table}

The sidebar now contains a clearly labelled \textbf{Department Pilots} route,
represented by the letter \textbf{P} when the rail is collapsed. It returns the
user to the department workspace containing Marketing, Sales, Finance,
Operations, Support and Human Resources.

\subsection{Marketing workspace consolidation}

Marketing previously existed in two disconnected locations: the Marketing Pilot
department and a separate Marketing Stream in the left sidebar. The Stream
contained the early manual content form, run history, approvals and calendar,
while the Pilot could produce an isolated artefact without launching the real
Marketing workflow.

Marketing now has one canonical home with the following components:

\begin{description}
    \item[AION Automation Studio:] The user describes a business outcome. AION
    reads the Boardroom assignment, Brand Foundation, tone of voice, target
    customer, offers and channel rules before launching the governed Marketing
    workflow.
    \item[Manual Creation Studio:] The user can directly specify the brief,
    objective, funnel goal, audience, persona, offer, channels, keywords,
    hashtags, hard rules, reference images, product assets and creative direction.
    \item[Shared workflow:] Automated and manual work use the same Brand
    Foundation, approval queue, content-generation path, run history and output
    store.
    \item[Operational navigation:] Stream, Calendar, Runs, Approvals and Settings
    remain accessible inside Marketing.
\end{description}

The obsolete standalone Marketing Stream link was removed. Historical routes are
redirected into Marketing, and \emph{Build with AION} now launches the real
Marketing workflow rather than producing an unconnected front-end artefact.

The consolidated foundation retains caption drafting, five-card carousel copy,
product and reference-image intake, the existing OpenAI image route, social-post
packaging, approvals, run history and failure evidence.

\subsection{Creative Intelligence Lab}

The Creative Intelligence Lab adds campaign-level creative orchestration above
the existing Marketing workflow. It converts one approved business idea into a
master concept, storyboard, generation jobs, platform renditions and a measurable
experiment.

\subsubsection{Creative planning}

For each campaign, the system records or derives:

\begin{itemize}
    \item campaign promise, objective, target audience and approved offer;
    \item tone, proof requirements and call to action;
    \item a five-scene structure covering hook, context, solution, proof and call
    to action;
    \item image, original-video, short-clip and carousel requirements;
    \item shot direction, captions, subtitle contracts and visual safe zones; and
    \item a stable plan identifier and SHA-256-based plan hash.
\end{itemize}

Generated claims must be supported by recorded business evidence. The engine is
not permitted to invent reviews, outcomes, prices or product capabilities.

\subsubsection{Platform renditions}

The platform-profile registry produces channel-specific rather than merely
resized outputs. Profiles currently cover Instagram feed, Reels and Stories;
Facebook feed and Reels; TikTok; YouTube Shorts; and LinkedIn feed and video.
The principal production contracts are:

\begin{itemize}
    \item $4{:}5$ portrait feed assets at $1080\times1350$ pixels;
    \item $9{:}16$ vertical video and clip assets at $1080\times1920$ pixels;
    \item platform-appropriate duration guidance;
    \item essential text, faces, logos and calls to action inside the central safe
    area; and
    \item burned-in caption guidance plus an editable subtitle-track contract for
    spoken video.
\end{itemize}

\subsubsection{Provider strategy}

Tessaris separates the creative operating system from the rendering model:

\begin{table}[htbp]
\centering
\small
\begin{tabular}{p{0.20\linewidth} p{0.28\linewidth} p{0.39\linewidth}}
\toprule
\textbf{Layer} & \textbf{Providers} & \textbf{Responsibility} \\
\midrule
Primary direct engines & OpenAI GPT Image/Sora; Google Imagen/Veo & Image and source-video generation through existing provider credentials \\
Optional specialists & Runway, Luma and fal & Premium or specialised models when directly connected and approved \\
Product references & Higgsfield and OpenArt & User-experience and capability references; not mandatory dependencies \\
Tessaris native layer & Tessaris/AION & Brand context, storyboards, prompts, platform adaptation, clipping, approvals, evidence and learning \\
\bottomrule
\end{tabular}
\caption{Provider-neutral creative-production architecture.}
\end{table}

This design prevents core Marketing capability from becoming locked to a
third-party wrapper whose models, prices or API availability may change.

\subsection{Controlled experimentation and winner learning}

The experiment governor creates a control and up to three challengers. Each
challenger changes exactly one variable:

\begin{itemize}
    \item the opening hook;
    \item the opening visual; or
    \item the call to action.
\end{itemize}

The default evidence contract requires at least 100 observations per variant, a
minimum relative improvement of 15\%, and no breach of complaint or truthfulness
guardrails. Challenger traffic is capped at 20\% until a result is accepted.

The evaluator can return four meaningful outcomes:

\begin{description}
    \item[Collect more evidence:] One or more variants have not reached the frozen
    sample threshold.
    \item[Keep control:] No challenger clears the improvement threshold, or an
    apparent improvement worsens a protected metric such as complaints.
    \item[Stop and review:] A false-claim, unsafe-content or equivalent guardrail
    is breached.
    \item[Winner candidate:] A challenger clears the primary threshold without
    breaching a guardrail. It still requires approval before traffic changes.
\end{description}

No evaluation automatically publishes content, purchases advertising, changes
campaign spend or promotes a winner. A live verification used simulated metrics
and correctly identified a challenger with a 30\% improvement as a candidate
while keeping automatic application disabled.

\subsection{Paid source-video rendering}

The paid-render layer implements persistent asynchronous jobs for OpenAI Sora and
Google Veo. A render job stores:

\begin{itemize}
    \item business workspace and campaign-plan reference;
    \item provider and exact model;
    \item full prompt and optional reference-asset path;
    \item aspect ratio, duration and resolution;
    \item rendition identifier and immutable approval hash;
    \item provider job identifier and provider response;
    \item completion status, error state and timestamps; and
    \item locally owned output metadata.
\end{itemize}

The lifecycle is deliberately separated into the following operations:

\begin{enumerate}
    \item \textbf{Prepare.} Build and persist the complete render contract without
    making a paid provider request.
    \item \textbf{Approve.} Bind the unchanged contract to an identified approver,
    an explicit paid-generation boolean and a positive maximum euro cost.
    \item \textbf{Submit.} Recalculate and compare the approval hash before calling
    the selected provider.
    \item \textbf{Poll.} Read the asynchronous provider status and fail closed on
    provider errors.
    \item \textbf{Acquire.} Download completed footage into the relevant Tessaris
    business workspace rather than relying on an expiring provider URL.
\end{enumerate}

A paid request without the explicit authority boolean returns HTTP status 403.
Verification prepared a dummy render contract and approved a maximum cost of
EUR~5, but deliberately did not submit the provider job; therefore no generation
charge was incurred.

Completed source footage is recorded with its local path, MIME type, byte count,
SHA-256 digest and an explicit \texttt{owned\_locally} flag. Provider credentials
are read at execution time and are never stored in render-job records.

\subsection{Native timeline, clipping and export}

Basic post-production is performed locally using the FFmpeg 7.1 binary already
available in the Tessaris Python environment. This avoids paying a premium
provider for deterministic editing tasks.

A timeline project contains an ordered set of source clips with start and end
times, labels and a target platform. Export performs the following operations:

\begin{itemize}
    \item validates every source path and trim boundary;
    \item trims and normalises each segment;
    \item scales while preserving aspect ratio;
    \item pads to the exact platform frame and produces square pixels;
    \item converts video to H.264 at 30 frames per second;
    \item preserves source audio as AAC stereo at 48~kHz;
    \item generates a silent audio stream when a source contains no audio;
    \item concatenates ordered segments and enables MP4 fast-start;
    \item creates a JPEG thumbnail; and
    \item stores output dimensions, byte count and SHA-256 evidence.
\end{itemize}

The desktop interface uses a minimal Electron inter-process bridge to open the
native macOS video picker. The renderer receives only the selected file path;
arbitrary filesystem browsing is not exposed to webpage code.

The completed visual editor extends this boundary with multi-select video,
image and audio pickers. Clips appear as draggable cards and retain independent
in/out points. The editor records bounded cuts, fades, dissolves, wipes and
slides; parses editable subtitle cues written as
\texttt{start-end|text}; applies timed text overlays; composites a local logo;
and loops and mixes background music at a user-selected volume. FFmpeg performs
all deterministic composition locally. Uploaded footage is not sent to a third
party by the editing operation.

Reference-image handling is provider aware. OpenAI jobs accept at most one
reference image, while Gemini jobs accept at most three. The selected paths form
part of the immutable approval hash. Gemini references are encoded into the
documented inline-data structure and reference-guided jobs are fixed to the
supported eight-second duration. No paid provider request occurs merely because
an image is selected.

Submitted render jobs are monitored by a daemon background poller. Completion
downloads the provider asset into the business workspace and records provenance,
byte count and SHA-256 evidence. A manual status button remains available, but
is no longer required for an unattended approved render to complete.

Physical verification produced a two-second MP4 with the following properties:

\begin{itemize}
    \item $1080\times1920$ pixels and $9{:}16$ display aspect ratio;
    \item H.264 High Profile video;
    \item AAC stereo audio at 48~kHz;
    \item 30 frames per second;
    \item locally stored thumbnail; and
    \item SHA-256 digest
    \texttt{80b1265eba012e690adb15cd833976f46e6dbca2d52e79dd0c7f2f5409a24b54}.
\end{itemize}

The verified output is stored at
\path{/Users/kevinrobinson/Tessaris/workspaces/demo/marketing/timelines/timeline_050a3f963e1545f2/instagram_reels_export.mp4}.

A second installed-application verification produced a 5.5-second two-clip
vertical export containing a fade transition, two burned subtitle cues, a timed
campaign overlay, a logo and mixed music. The service also proved that an
eight-second default trim is clamped safely to each source file's real duration.
The final visual-QA asset has SHA-256 digest
\texttt{9e8d813556c4af1198bbd416a5ef4a09bd0d51d38e7ef3de39ac54b069fb5804}.

\subsection{Premium motion, narration and publication authority}

Each clip now carries an optional static, slow zoom-in, slow zoom-out, left-pan
or right-pan motion treatment. Motion is calculated against the real source
duration and rendered locally. Titles support no animation, fade or slide-up,
and the user may select clean editorial, bold commercial or cinematic brand
treatment.

The audio graph accepts an independent voice-over track with start-time and gain
controls. When enabled, side-chain compression lowers the background music under
narration. Music receives bounded fades and the complete mix can be normalised to
the social or broadcast loudness preset. No cloud service is required for these
operations.

The desktop can also record narration directly from the authorised microphone.
The browser recorder sends the resulting Opus/WebM asset only to the local
Tessaris backend, which stores it under the current business workspace with byte
count and SHA-256 evidence before it enters the same mastering graph.

A governed publication package freezes the exact finished-asset digest, caption,
call to action, channel, organic or paid mode, schedule and budget. Exact approval
is recorded separately. The package reports missing Meta, Google Ads or platform
credentials and has no external execution endpoint in this increment, so neither
publication nor advertising spend can occur accidentally.

Campaign evidence records impressions, clicks, qualified actions, conversions,
spend, revenue, complaints and false or unsafe claims. Tessaris derives click
rate, qualified-action rate, conversion rate, cost per qualified action and
return on advertising spend. Unsafe evidence changes the package state to an
escalated measurement rather than a winner.

The premium physical proof produced a 5.4-second $1080\times1920$ H.264/AAC
master containing two motion-treated clips, dissolve, animated title, subtitles,
logo, music, separate narration, automatic ducking and social loudness
mastering. Its SHA-256 digest is
\texttt{816264199101da6915bb50dd192b963bc174e2ba755ff30766171aa4f85d5e4b}.

\subsection{HTTP interfaces}

The provider-neutral planning and experimentation boundary exposes:

\begin{itemize}
    \item \path{GET /api/aion/marketing/creative/capabilities}
    \item \path{GET /api/aion/marketing/creative/platforms}
    \item \path{POST /api/aion/marketing/creative/plan}
    \item \path{POST /api/aion/marketing/creative/experiment/evaluate}
\end{itemize}

The paid-render and local-timeline boundary exposes:

\begin{itemize}
    \item \path{GET /api/aion/marketing/creative/{workspace_id}/renders}
    \item \path{POST /api/aion/marketing/creative/{workspace_id}/renders}
    \item \path{POST /api/aion/marketing/creative/{workspace_id}/renders/{job_id}/approve}
    \item \path{POST /api/aion/marketing/creative/{workspace_id}/renders/{job_id}/submit}
    \item \path{POST /api/aion/marketing/creative/{workspace_id}/renders/{job_id}/poll}
    \item \path{POST /api/aion/marketing/creative/{workspace_id}/timelines}
    \item \path{POST /api/aion/marketing/creative/{workspace_id}/timelines/{timeline_id}/export}
\end{itemize}

The routers are mounted in both the historical full backend and the lean desktop
backend. The latter is the runtime used by the installed macOS application.

\subsection{Principal implementation files}

\begin{itemize}
    \item Desktop interface and department routing:
    \path{desktop/mac/src/app.js}
    \item Native video-picker boundary:
    \path{desktop/mac/electron/main.js} and
    \path{desktop/mac/electron/preload.js}
    \item Creative HTTP API:
    \path{backend/modules/aion_business/api/marketing_creative_api.py}
    \item Campaign and experiment engine:
    \path{backend/modules/aion_business/runtime/marketing_creative_engine.py}
    \item Provider request contracts:
    \path{backend/modules/aion_business/runtime/creative_provider_adapters.py}
    \item Persistent paid-render service:
    \path{backend/modules/aion_business/runtime/creative_render_service.py}
    \item Automatic provider completion monitor:
    \path{backend/modules/aion_business/runtime/creative_render_polling_service.py}
    \item Local timeline and export service:
    \path{backend/modules/aion_business/runtime/creative_timeline_service.py}
    \item Publication authority and campaign evidence service:
    \path{backend/modules/aion_business/runtime/creative_publication_service.py}
    \item Lean desktop application boundary:
    \path{backend/desktop_app.py}
\end{itemize}

\subsection{Verification and recovery}

The department-shell standardisation passed 37 relevant regression checks. The
Marketing consolidation passed 65 focused checks. Additional focused checks
cover creative-plan construction, provider request contracts, evidence thresholds,
truthfulness safeguards, paid-render approval and timeline export.

The installed desktop backend reports healthy in lean
\texttt{tessaris\_desktop} mode. Relevant recoverable application snapshots are:

\begin{itemize}
    \item \path{/Applications/Tessaris.app.backup-20260810-before-support-layout-consolidation}
    \item \path{/Applications/Tessaris.app.backup-20260810-before-unified-department-shells}
    \item \path{/Applications/Tessaris.app.backup-20260810-before-unified-marketing-studio}
    \item \path{/Applications/Tessaris.app.backup-20260810-before-creative-intelligence-lab}
    \item \path{/Applications/Tessaris.app.backup-20260810-before-paid-render-and-local-timeline}
    \item \path{/Applications/Tessaris.app.backup-20260810-before-department-pilots-sidebar-link}
    \item \path{/Applications/Tessaris.app.backup-20260810-before-marketing-editor-completion}
    \item \path{/Applications/Tessaris.app.backup-20260810-before-premium-marketing-production}
    \item \path{/Applications/Tessaris.app.backup-20260810-before-local-narration-recorder}
\end{itemize}

No Git commit was made during these increments.

\subsection{Honest remaining boundary}

The architecture, governed render execution and physical platform export are now
operational. The remaining work is primarily production-channel completion and a
richer editor rather than another foundational redesign:

\begin{itemize}
    \item free-form keyframe curves beyond the safe motion presets and fine-grained caption typography;
    \item provider-backed speech generation as an optional alternative to imported or locally recorded narration;
    \item provider-native webhooks as an optimisation over the operational automatic poller;
    \item optional Runway, Luma and fal credential/execution adapters;
    \item executable Meta, Google and other publishing or advertising connectors after business credentials are supplied;
    \item automatic conversion feedback from Sales and imported platform analytics; and
    \item real campaign validation across multiple businesses and channels.
\end{itemize}

The intended operating rule remains: Tessaris may plan, generate, edit and learn
broadly, but paid generation, publishing, advertising spend and winner promotion
remain bounded by explicit business authority and recorded evidence.


\section{Unified Department Pilots and Marketing Creative Production Stack}
\label{sec:tessaris-department-pilots-marketing-creative}

\subsection{Scope and design objective}

This development phase consolidated the Tessaris department workspaces and
advanced Marketing from an early social-post prototype into a governed creative
production and experimentation system. The work covered five connected layers:

\begin{enumerate}
    \item a consistent department page structure for Marketing, Sales, Finance,
    Operations, Support and Human Resources;
    \item department-scoped AION/Pilot conversations and Boardroom assignment
    packages;
    \item consolidation of the former standalone Marketing Stream into the
    Marketing Pilot workspace;
    \item a provider-neutral Creative Intelligence Lab for images, video, clips,
    platform renditions and controlled experiments; and
    \item exact-approved paid video rendering plus locally owned timeline and
    clip export.
\end{enumerate}

The central architectural principle is that Tessaris owns the business context,
creative direction, governance, asset provenance and learning loop. External
generative-AI providers remain replaceable production engines. This allows the
product to perform as much work as possible locally while using specialist or
premium services only where model-scale generation is genuinely required.

\subsection{Unified department workspace shell}

The earlier department pages had diverged. Some routes rendered a full-screen
spatial Boardroom, some retained an obsolete black Pilot terminal, and some hid
the newer operational workspace. Support and Human Resources were especially
affected because legacy containers could overlap or replace the newer interfaces.

Every principal department now follows the same hierarchy:

\begin{enumerate}
    \item \textbf{Compact department director.} A small department-specific AION
    character identifies the function and its advisory boundary. It replaces the
    oversized spatial room on operational pages.
    \item \textbf{Department-scoped conversation.} The \emph{Talk to Pilot}
    control opens a conversation whose context is limited to the active
    department. For example, the Support Pilot discusses customer cases,
    policies and escalations rather than acting as a generic assistant.
    \item \textbf{Boardroom assignment package.} The package is retained because
    it contains the Boardroom meeting instructions, goals, discovery requirements,
    linked goal sheet and expected outputs. The \emph{Open Package} control is
    available consistently.
    \item \textbf{Operational workspace.} The real department system appears
    beneath the package: Sales Revenue, Finance, Support Case Centre, HR people
    and permissions, Marketing, or the Operations foundation.
\end{enumerate}

Legacy black Pilot containers and oversized spatial canvases are explicitly
excluded from these routes when they would override the operational workspace.
The installed department outcomes are:

\begin{table}[htbp]
\centering
\small
\begin{tabular}{p{0.17\linewidth} p{0.70\linewidth}}
\toprule
\textbf{Department} & \textbf{Mounted operational workspace} \\
\midrule
Marketing & Unified automation, creative production, approvals, calendar and history \\
Sales & Canonical CRM, governed communications and Sales Revenue workspace \\
Finance & Finance operating workspace, evidence, bookkeeping and reporting controls \\
Operations & Unified department shell and foundation for the Operations build \\
Support & Canonical Support Case Centre, policy authority and escalation workflow \\
Human Resources & Organisation, people, roles, permissions and authority workspace \\
\bottomrule
\end{tabular}
\caption{Standardised department workspace structure.}
\end{table}

The sidebar now contains a clearly labelled \textbf{Department Pilots} route,
represented by the letter \textbf{P} when the rail is collapsed. It returns the
user to the department workspace containing Marketing, Sales, Finance,
Operations, Support and Human Resources.

\subsection{Marketing workspace consolidation}

Marketing previously existed in two disconnected locations: the Marketing Pilot
department and a separate Marketing Stream in the left sidebar. The Stream
contained the early manual content form, run history, approvals and calendar,
while the Pilot could produce an isolated artefact without launching the real
Marketing workflow.

Marketing now has one canonical home with the following components:

\begin{description}
    \item[AION Automation Studio:] The user describes a business outcome. AION
    reads the Boardroom assignment, Brand Foundation, tone of voice, target
    customer, offers and channel rules before launching the governed Marketing
    workflow.
    \item[Manual Creation Studio:] The user can directly specify the brief,
    objective, funnel goal, audience, persona, offer, channels, keywords,
    hashtags, hard rules, reference images, product assets and creative direction.
    \item[Shared workflow:] Automated and manual work use the same Brand
    Foundation, approval queue, content-generation path, run history and output
    store.
    \item[Operational navigation:] Stream, Calendar, Runs, Approvals and Settings
    remain accessible inside Marketing.
\end{description}

The obsolete standalone Marketing Stream link was removed. Historical routes are
redirected into Marketing, and \emph{Build with AION} now launches the real
Marketing workflow rather than producing an unconnected front-end artefact.

The consolidated foundation retains caption drafting, five-card carousel copy,
product and reference-image intake, the existing OpenAI image route, social-post
packaging, approvals, run history and failure evidence.

\subsection{Creative Intelligence Lab}

The Creative Intelligence Lab adds campaign-level creative orchestration above
the existing Marketing workflow. It converts one approved business idea into a
master concept, storyboard, generation jobs, platform renditions and a measurable
experiment.

\subsubsection{Creative planning}

For each campaign, the system records or derives:

\begin{itemize}
    \item campaign promise, objective, target audience and approved offer;
    \item tone, proof requirements and call to action;
    \item a five-scene structure covering hook, context, solution, proof and call
    to action;
    \item image, original-video, short-clip and carousel requirements;
    \item shot direction, captions, subtitle contracts and visual safe zones; and
    \item a stable plan identifier and SHA-256-based plan hash.
\end{itemize}

Generated claims must be supported by recorded business evidence. The engine is
not permitted to invent reviews, outcomes, prices or product capabilities.

\subsubsection{Platform renditions}

The platform-profile registry produces channel-specific rather than merely
resized outputs. Profiles currently cover Instagram feed, Reels and Stories;
Facebook feed and Reels; TikTok; YouTube Shorts; and LinkedIn feed and video.
The principal production contracts are:

\begin{itemize}
    \item $4{:}5$ portrait feed assets at $1080\times1350$ pixels;
    \item $9{:}16$ vertical video and clip assets at $1080\times1920$ pixels;
    \item platform-appropriate duration guidance;
    \item essential text, faces, logos and calls to action inside the central safe
    area; and
    \item burned-in caption guidance plus an editable subtitle-track contract for
    spoken video.
\end{itemize}

\subsubsection{Provider strategy}

Tessaris separates the creative operating system from the rendering model:

\begin{table}[htbp]
\centering
\small
\begin{tabular}{p{0.20\linewidth} p{0.28\linewidth} p{0.39\linewidth}}
\toprule
\textbf{Layer} & \textbf{Providers} & \textbf{Responsibility} \\
\midrule
Primary direct engines & OpenAI GPT Image/Sora; Google Imagen/Veo & Image and source-video generation through existing provider credentials \\
Optional specialists & Runway, Luma and fal & Premium or specialised models when directly connected and approved \\
Product references & Higgsfield and OpenArt & User-experience and capability references; not mandatory dependencies \\
Tessaris native layer & Tessaris/AION & Brand context, storyboards, prompts, platform adaptation, clipping, approvals, evidence and learning \\
\bottomrule
\end{tabular}
\caption{Provider-neutral creative-production architecture.}
\end{table}

This design prevents core Marketing capability from becoming locked to a
third-party wrapper whose models, prices or API availability may change.

\subsection{Controlled experimentation and winner learning}

The experiment governor creates a control and up to three challengers. Each
challenger changes exactly one variable:

\begin{itemize}
    \item the opening hook;
    \item the opening visual; or
    \item the call to action.
\end{itemize}

The default evidence contract requires at least 100 observations per variant, a
minimum relative improvement of 15\%, and no breach of complaint or truthfulness
guardrails. Challenger traffic is capped at 20\% until a result is accepted.

The evaluator can return four meaningful outcomes:

\begin{description}
    \item[Collect more evidence:] One or more variants have not reached the frozen
    sample threshold.
    \item[Keep control:] No challenger clears the improvement threshold, or an
    apparent improvement worsens a protected metric such as complaints.
    \item[Stop and review:] A false-claim, unsafe-content or equivalent guardrail
    is breached.
    \item[Winner candidate:] A challenger clears the primary threshold without
    breaching a guardrail. It still requires approval before traffic changes.
\end{description}

No evaluation automatically publishes content, purchases advertising, changes
campaign spend or promotes a winner. A live verification used simulated metrics
and correctly identified a challenger with a 30\% improvement as a candidate
while keeping automatic application disabled.

\subsection{Paid source-video rendering}

The paid-render layer implements persistent asynchronous jobs for OpenAI Sora and
Google Veo. A render job stores:

\begin{itemize}
    \item business workspace and campaign-plan reference;
    \item provider and exact model;
    \item full prompt and optional reference-asset path;
    \item aspect ratio, duration and resolution;
    \item rendition identifier and immutable approval hash;
    \item provider job identifier and provider response;
    \item completion status, error state and timestamps; and
    \item locally owned output metadata.
\end{itemize}

The lifecycle is deliberately separated into the following operations:

\begin{enumerate}
    \item \textbf{Prepare.} Build and persist the complete render contract without
    making a paid provider request.
    \item \textbf{Approve.} Bind the unchanged contract to an identified approver,
    an explicit paid-generation boolean and a positive maximum euro cost.
    \item \textbf{Submit.} Recalculate and compare the approval hash before calling
    the selected provider.
    \item \textbf{Poll.} Read the asynchronous provider status and fail closed on
    provider errors.
    \item \textbf{Acquire.} Download completed footage into the relevant Tessaris
    business workspace rather than relying on an expiring provider URL.
\end{enumerate}

A paid request without the explicit authority boolean returns HTTP status 403.
Verification prepared a dummy render contract and approved a maximum cost of
EUR~5, but deliberately did not submit the provider job; therefore no generation
charge was incurred.

Completed source footage is recorded with its local path, MIME type, byte count,
SHA-256 digest and an explicit \texttt{owned\_locally} flag. Provider credentials
are read at execution time and are never stored in render-job records.

\subsection{Native timeline, clipping and export}

Basic post-production is performed locally using the FFmpeg 7.1 binary already
available in the Tessaris Python environment. This avoids paying a premium
provider for deterministic editing tasks.

A timeline project contains an ordered set of source clips with start and end
times, labels and a target platform. Export performs the following operations:

\begin{itemize}
    \item validates every source path and trim boundary;
    \item trims and normalises each segment;
    \item scales while preserving aspect ratio;
    \item pads to the exact platform frame and produces square pixels;
    \item converts video to H.264 at 30 frames per second;
    \item preserves source audio as AAC stereo at 48~kHz;
    \item generates a silent audio stream when a source contains no audio;
    \item concatenates ordered segments and enables MP4 fast-start;
    \item creates a JPEG thumbnail; and
    \item stores output dimensions, byte count and SHA-256 evidence.
\end{itemize}

The desktop interface uses a minimal Electron inter-process bridge to open the
native macOS video picker. The renderer receives only the selected file path;
arbitrary filesystem browsing is not exposed to webpage code.

The completed visual editor extends this boundary with multi-select video,
image and audio pickers. Clips appear as draggable cards and retain independent
in/out points. The editor records bounded cuts, fades, dissolves, wipes and
slides; parses editable subtitle cues written as
\texttt{start-end|text}; applies timed text overlays; composites a local logo;
and loops and mixes background music at a user-selected volume. FFmpeg performs
all deterministic composition locally. Uploaded footage is not sent to a third
party by the editing operation.

Reference-image handling is provider aware. OpenAI jobs accept at most one
reference image, while Gemini jobs accept at most three. The selected paths form
part of the immutable approval hash. Gemini references are encoded into the
documented inline-data structure and reference-guided jobs are fixed to the
supported eight-second duration. No paid provider request occurs merely because
an image is selected.

Submitted render jobs are monitored by a daemon background poller. Completion
downloads the provider asset into the business workspace and records provenance,
byte count and SHA-256 evidence. A manual status button remains available, but
is no longer required for an unattended approved render to complete.

Physical verification produced a two-second MP4 with the following properties:

\begin{itemize}
    \item $1080\times1920$ pixels and $9{:}16$ display aspect ratio;
    \item H.264 High Profile video;
    \item AAC stereo audio at 48~kHz;
    \item 30 frames per second;
    \item locally stored thumbnail; and
    \item SHA-256 digest
    \texttt{80b1265eba012e690adb15cd833976f46e6dbca2d52e79dd0c7f2f5409a24b54}.
\end{itemize}

The verified output is stored at
\path{/Users/kevinrobinson/Tessaris/workspaces/demo/marketing/timelines/timeline_050a3f963e1545f2/instagram_reels_export.mp4}.

A second installed-application verification produced a 5.5-second two-clip
vertical export containing a fade transition, two burned subtitle cues, a timed
campaign overlay, a logo and mixed music. The service also proved that an
eight-second default trim is clamped safely to each source file's real duration.
The final visual-QA asset has SHA-256 digest
\texttt{9e8d813556c4af1198bbd416a5ef4a09bd0d51d38e7ef3de39ac54b069fb5804}.

\subsection{Premium motion, narration and publication authority}

Each clip now carries an optional static, slow zoom-in, slow zoom-out, left-pan
or right-pan motion treatment. Motion is calculated against the real source
duration and rendered locally. Titles support no animation, fade or slide-up,
and the user may select clean editorial, bold commercial or cinematic brand
treatment.

The audio graph accepts an independent voice-over track with start-time and gain
controls. When enabled, side-chain compression lowers the background music under
narration. Music receives bounded fades and the complete mix can be normalised to
the social or broadcast loudness preset. No cloud service is required for these
operations.

The desktop can also record narration directly from the authorised microphone.
The browser recorder sends the resulting Opus/WebM asset only to the local
Tessaris backend, which stores it under the current business workspace with byte
count and SHA-256 evidence before it enters the same mastering graph.

A governed publication package freezes the exact finished-asset digest, caption,
call to action, channel, organic or paid mode, schedule and budget. Exact approval
is recorded separately. The package reports missing Meta, Google Ads or platform
credentials and has no external execution endpoint in this increment, so neither
publication nor advertising spend can occur accidentally.

Campaign evidence records impressions, clicks, qualified actions, conversions,
spend, revenue, complaints and false or unsafe claims. Tessaris derives click
rate, qualified-action rate, conversion rate, cost per qualified action and
return on advertising spend. Unsafe evidence changes the package state to an
escalated measurement rather than a winner.

The premium physical proof produced a 5.4-second $1080\times1920$ H.264/AAC
master containing two motion-treated clips, dissolve, animated title, subtitles,
logo, music, separate narration, automatic ducking and social loudness
mastering. Its SHA-256 digest is
\texttt{816264199101da6915bb50dd192b963bc174e2ba755ff30766171aa4f85d5e4b}.

\subsection{HTTP interfaces}

The provider-neutral planning and experimentation boundary exposes:

\begin{itemize}
    \item \path{GET /api/aion/marketing/creative/capabilities}
    \item \path{GET /api/aion/marketing/creative/platforms}
    \item \path{POST /api/aion/marketing/creative/plan}
    \item \path{POST /api/aion/marketing/creative/experiment/evaluate}
\end{itemize}

The paid-render and local-timeline boundary exposes:

\begin{itemize}
    \item \path{GET /api/aion/marketing/creative/{workspace_id}/renders}
    \item \path{POST /api/aion/marketing/creative/{workspace_id}/renders}
    \item \path{POST /api/aion/marketing/creative/{workspace_id}/renders/{job_id}/approve}
    \item \path{POST /api/aion/marketing/creative/{workspace_id}/renders/{job_id}/submit}
    \item \path{POST /api/aion/marketing/creative/{workspace_id}/renders/{job_id}/poll}
    \item \path{POST /api/aion/marketing/creative/{workspace_id}/timelines}
    \item \path{POST /api/aion/marketing/creative/{workspace_id}/timelines/{timeline_id}/export}
\end{itemize}

The routers are mounted in both the historical full backend and the lean desktop
backend. The latter is the runtime used by the installed macOS application.

\subsection{Principal implementation files}

\begin{itemize}
    \item Desktop interface and department routing:
    \path{desktop/mac/src/app.js}
    \item Native video-picker boundary:
    \path{desktop/mac/electron/main.js} and
    \path{desktop/mac/electron/preload.js}
    \item Creative HTTP API:
    \path{backend/modules/aion_business/api/marketing_creative_api.py}
    \item Campaign and experiment engine:
    \path{backend/modules/aion_business/runtime/marketing_creative_engine.py}
    \item Provider request contracts:
    \path{backend/modules/aion_business/runtime/creative_provider_adapters.py}
    \item Persistent paid-render service:
    \path{backend/modules/aion_business/runtime/creative_render_service.py}
    \item Automatic provider completion monitor:
    \path{backend/modules/aion_business/runtime/creative_render_polling_service.py}
    \item Local timeline and export service:
    \path{backend/modules/aion_business/runtime/creative_timeline_service.py}
    \item Publication authority and campaign evidence service:
    \path{backend/modules/aion_business/runtime/creative_publication_service.py}
    \item Lean desktop application boundary:
    \path{backend/desktop_app.py}
\end{itemize}

\subsection{Verification and recovery}

The department-shell standardisation passed 37 relevant regression checks. The
Marketing consolidation passed 65 focused checks. Additional focused checks
cover creative-plan construction, provider request contracts, evidence thresholds,
truthfulness safeguards, paid-render approval and timeline export.

The installed desktop backend reports healthy in lean
\texttt{tessaris\_desktop} mode. Relevant recoverable application snapshots are:

\begin{itemize}
    \item \path{/Applications/Tessaris.app.backup-20260810-before-support-layout-consolidation}
    \item \path{/Applications/Tessaris.app.backup-20260810-before-unified-department-shells}
    \item \path{/Applications/Tessaris.app.backup-20260810-before-unified-marketing-studio}
    \item \path{/Applications/Tessaris.app.backup-20260810-before-creative-intelligence-lab}
    \item \path{/Applications/Tessaris.app.backup-20260810-before-paid-render-and-local-timeline}
    \item \path{/Applications/Tessaris.app.backup-20260810-before-department-pilots-sidebar-link}
    \item \path{/Applications/Tessaris.app.backup-20260810-before-marketing-editor-completion}
    \item \path{/Applications/Tessaris.app.backup-20260810-before-premium-marketing-production}
    \item \path{/Applications/Tessaris.app.backup-20260810-before-local-narration-recorder}
\end{itemize}

No Git commit was made during these increments.

\subsection{Honest remaining boundary}

The architecture, governed render execution and physical platform export are now
operational. The remaining work is primarily production-channel completion and a
richer editor rather than another foundational redesign:

\begin{itemize}
    \item free-form keyframe curves beyond the safe motion presets and fine-grained caption typography;
    \item provider-backed speech generation as an optional alternative to imported or locally recorded narration;
    \item provider-native webhooks as an optimisation over the operational automatic poller;
    \item optional Runway, Luma and fal credential/execution adapters;
    \item executable Meta, Google and other publishing or advertising connectors after business credentials are supplied;
    \item automatic conversion feedback from Sales and imported platform analytics; and
    \item real campaign validation across multiple businesses and channels.
\end{itemize}

The intended operating rule remains: Tessaris may plan, generate, edit and learn
broadly, but paid generation, publishing, advertising spend and winner promotion
remain bounded by explicit business authority and recorded evidence.


\section{Unified Department Pilots and Marketing Creative Production Stack}
\label{sec:tessaris-department-pilots-marketing-creative}

\subsection{Scope and design objective}

This development phase consolidated the Tessaris department workspaces and
advanced Marketing from an early social-post prototype into a governed creative
production and experimentation system. The work covered five connected layers:

\begin{enumerate}
    \item a consistent department page structure for Marketing, Sales, Finance,
    Operations, Support and Human Resources;
    \item department-scoped AION/Pilot conversations and Boardroom assignment
    packages;
    \item consolidation of the former standalone Marketing Stream into the
    Marketing Pilot workspace;
    \item a provider-neutral Creative Intelligence Lab for images, video, clips,
    platform renditions and controlled experiments; and
    \item exact-approved paid video rendering plus locally owned timeline and
    clip export.
\end{enumerate}

The central architectural principle is that Tessaris owns the business context,
creative direction, governance, asset provenance and learning loop. External
generative-AI providers remain replaceable production engines. This allows the
product to perform as much work as possible locally while using specialist or
premium services only where model-scale generation is genuinely required.

\subsection{Unified department workspace shell}

The earlier department pages had diverged. Some routes rendered a full-screen
spatial Boardroom, some retained an obsolete black Pilot terminal, and some hid
the newer operational workspace. Support and Human Resources were especially
affected because legacy containers could overlap or replace the newer interfaces.

Every principal department now follows the same hierarchy:

\begin{enumerate}
    \item \textbf{Compact department director.} A small department-specific AION
    character identifies the function and its advisory boundary. It replaces the
    oversized spatial room on operational pages.
    \item \textbf{Department-scoped conversation.} The \emph{Talk to Pilot}
    control opens a conversation whose context is limited to the active
    department. For example, the Support Pilot discusses customer cases,
    policies and escalations rather than acting as a generic assistant.
    \item \textbf{Boardroom assignment package.} The package is retained because
    it contains the Boardroom meeting instructions, goals, discovery requirements,
    linked goal sheet and expected outputs. The \emph{Open Package} control is
    available consistently.
    \item \textbf{Operational workspace.} The real department system appears
    beneath the package: Sales Revenue, Finance, Support Case Centre, HR people
    and permissions, Marketing, or the Operations foundation.
\end{enumerate}

Legacy black Pilot containers and oversized spatial canvases are explicitly
excluded from these routes when they would override the operational workspace.
The installed department outcomes are:

\begin{table}[htbp]
\centering
\small
\begin{tabular}{p{0.17\linewidth} p{0.70\linewidth}}
\toprule
\textbf{Department} & \textbf{Mounted operational workspace} \\
\midrule
Marketing & Unified automation, creative production, approvals, calendar and history \\
Sales & Canonical CRM, governed communications and Sales Revenue workspace \\
Finance & Finance operating workspace, evidence, bookkeeping and reporting controls \\
Operations & Unified department shell and foundation for the Operations build \\
Support & Canonical Support Case Centre, policy authority and escalation workflow \\
Human Resources & Organisation, people, roles, permissions and authority workspace \\
\bottomrule
\end{tabular}
\caption{Standardised department workspace structure.}
\end{table}

The sidebar now contains a clearly labelled \textbf{Department Pilots} route,
represented by the letter \textbf{P} when the rail is collapsed. It returns the
user to the department workspace containing Marketing, Sales, Finance,
Operations, Support and Human Resources.

\subsection{Marketing workspace consolidation}

Marketing previously existed in two disconnected locations: the Marketing Pilot
department and a separate Marketing Stream in the left sidebar. The Stream
contained the early manual content form, run history, approvals and calendar,
while the Pilot could produce an isolated artefact without launching the real
Marketing workflow.

Marketing now has one canonical home with the following components:

\begin{description}
    \item[AION Automation Studio:] The user describes a business outcome. AION
    reads the Boardroom assignment, Brand Foundation, tone of voice, target
    customer, offers and channel rules before launching the governed Marketing
    workflow.
    \item[Manual Creation Studio:] The user can directly specify the brief,
    objective, funnel goal, audience, persona, offer, channels, keywords,
    hashtags, hard rules, reference images, product assets and creative direction.
    \item[Shared workflow:] Automated and manual work use the same Brand
    Foundation, approval queue, content-generation path, run history and output
    store.
    \item[Operational navigation:] Stream, Calendar, Runs, Approvals and Settings
    remain accessible inside Marketing.
\end{description}

The obsolete standalone Marketing Stream link was removed. Historical routes are
redirected into Marketing, and \emph{Build with AION} now launches the real
Marketing workflow rather than producing an unconnected front-end artefact.

The consolidated foundation retains caption drafting, five-card carousel copy,
product and reference-image intake, the existing OpenAI image route, social-post
packaging, approvals, run history and failure evidence.

\subsection{Creative Intelligence Lab}

The Creative Intelligence Lab adds campaign-level creative orchestration above
the existing Marketing workflow. It converts one approved business idea into a
master concept, storyboard, generation jobs, platform renditions and a measurable
experiment.

\subsubsection{Creative planning}

For each campaign, the system records or derives:

\begin{itemize}
    \item campaign promise, objective, target audience and approved offer;
    \item tone, proof requirements and call to action;
    \item a five-scene structure covering hook, context, solution, proof and call
    to action;
    \item image, original-video, short-clip and carousel requirements;
    \item shot direction, captions, subtitle contracts and visual safe zones; and
    \item a stable plan identifier and SHA-256-based plan hash.
\end{itemize}

Generated claims must be supported by recorded business evidence. The engine is
not permitted to invent reviews, outcomes, prices or product capabilities.

\subsubsection{Platform renditions}

The platform-profile registry produces channel-specific rather than merely
resized outputs. Profiles currently cover Instagram feed, Reels and Stories;
Facebook feed and Reels; TikTok; YouTube Shorts; and LinkedIn feed and video.
The principal production contracts are:

\begin{itemize}
    \item $4{:}5$ portrait feed assets at $1080\times1350$ pixels;
    \item $9{:}16$ vertical video and clip assets at $1080\times1920$ pixels;
    \item platform-appropriate duration guidance;
    \item essential text, faces, logos and calls to action inside the central safe
    area; and
    \item burned-in caption guidance plus an editable subtitle-track contract for
    spoken video.
\end{itemize}

\subsubsection{Provider strategy}

Tessaris separates the creative operating system from the rendering model:

\begin{table}[htbp]
\centering
\small
\begin{tabular}{p{0.20\linewidth} p{0.28\linewidth} p{0.39\linewidth}}
\toprule
\textbf{Layer} & \textbf{Providers} & \textbf{Responsibility} \\
\midrule
Primary direct engines & OpenAI GPT Image/Sora; Google Imagen/Veo & Image and source-video generation through existing provider credentials \\
Optional specialists & Runway, Luma and fal & Premium or specialised models when directly connected and approved \\
Product references & Higgsfield and OpenArt & User-experience and capability references; not mandatory dependencies \\
Tessaris native layer & Tessaris/AION & Brand context, storyboards, prompts, platform adaptation, clipping, approvals, evidence and learning \\
\bottomrule
\end{tabular}
\caption{Provider-neutral creative-production architecture.}
\end{table}

This design prevents core Marketing capability from becoming locked to a
third-party wrapper whose models, prices or API availability may change.

\subsection{Controlled experimentation and winner learning}

The experiment governor creates a control and up to three challengers. Each
challenger changes exactly one variable:

\begin{itemize}
    \item the opening hook;
    \item the opening visual; or
    \item the call to action.
\end{itemize}

The default evidence contract requires at least 100 observations per variant, a
minimum relative improvement of 15\%, and no breach of complaint or truthfulness
guardrails. Challenger traffic is capped at 20\% until a result is accepted.

The evaluator can return four meaningful outcomes:

\begin{description}
    \item[Collect more evidence:] One or more variants have not reached the frozen
    sample threshold.
    \item[Keep control:] No challenger clears the improvement threshold, or an
    apparent improvement worsens a protected metric such as complaints.
    \item[Stop and review:] A false-claim, unsafe-content or equivalent guardrail
    is breached.
    \item[Winner candidate:] A challenger clears the primary threshold without
    breaching a guardrail. It still requires approval before traffic changes.
\end{description}

No evaluation automatically publishes content, purchases advertising, changes
campaign spend or promotes a winner. A live verification used simulated metrics
and correctly identified a challenger with a 30\% improvement as a candidate
while keeping automatic application disabled.

\subsection{Paid source-video rendering}

The paid-render layer implements persistent asynchronous jobs for OpenAI Sora and
Google Veo. A render job stores:

\begin{itemize}
    \item business workspace and campaign-plan reference;
    \item provider and exact model;
    \item full prompt and optional reference-asset path;
    \item aspect ratio, duration and resolution;
    \item rendition identifier and immutable approval hash;
    \item provider job identifier and provider response;
    \item completion status, error state and timestamps; and
    \item locally owned output metadata.
\end{itemize}

The lifecycle is deliberately separated into the following operations:

\begin{enumerate}
    \item \textbf{Prepare.} Build and persist the complete render contract without
    making a paid provider request.
    \item \textbf{Approve.} Bind the unchanged contract to an identified approver,
    an explicit paid-generation boolean and a positive maximum euro cost.
    \item \textbf{Submit.} Recalculate and compare the approval hash before calling
    the selected provider.
    \item \textbf{Poll.} Read the asynchronous provider status and fail closed on
    provider errors.
    \item \textbf{Acquire.} Download completed footage into the relevant Tessaris
    business workspace rather than relying on an expiring provider URL.
\end{enumerate}

A paid request without the explicit authority boolean returns HTTP status 403.
Verification prepared a dummy render contract and approved a maximum cost of
EUR~5, but deliberately did not submit the provider job; therefore no generation
charge was incurred.

Completed source footage is recorded with its local path, MIME type, byte count,
SHA-256 digest and an explicit \texttt{owned\_locally} flag. Provider credentials
are read at execution time and are never stored in render-job records.

\subsection{Native timeline, clipping and export}

Basic post-production is performed locally using the FFmpeg 7.1 binary already
available in the Tessaris Python environment. This avoids paying a premium
provider for deterministic editing tasks.

A timeline project contains an ordered set of source clips with start and end
times, labels and a target platform. Export performs the following operations:

\begin{itemize}
    \item validates every source path and trim boundary;
    \item trims and normalises each segment;
    \item scales while preserving aspect ratio;
    \item pads to the exact platform frame and produces square pixels;
    \item converts video to H.264 at 30 frames per second;
    \item preserves source audio as AAC stereo at 48~kHz;
    \item generates a silent audio stream when a source contains no audio;
    \item concatenates ordered segments and enables MP4 fast-start;
    \item creates a JPEG thumbnail; and
    \item stores output dimensions, byte count and SHA-256 evidence.
\end{itemize}

The desktop interface uses a minimal Electron inter-process bridge to open the
native macOS video picker. The renderer receives only the selected file path;
arbitrary filesystem browsing is not exposed to webpage code.

The completed visual editor extends this boundary with multi-select video,
image and audio pickers. Clips appear as draggable cards and retain independent
in/out points. The editor records bounded cuts, fades, dissolves, wipes and
slides; parses editable subtitle cues written as
\texttt{start-end|text}; applies timed text overlays; composites a local logo;
and loops and mixes background music at a user-selected volume. FFmpeg performs
all deterministic composition locally. Uploaded footage is not sent to a third
party by the editing operation.

Reference-image handling is provider aware. OpenAI jobs accept at most one
reference image, while Gemini jobs accept at most three. The selected paths form
part of the immutable approval hash. Gemini references are encoded into the
documented inline-data structure and reference-guided jobs are fixed to the
supported eight-second duration. No paid provider request occurs merely because
an image is selected.

Submitted render jobs are monitored by a daemon background poller. Completion
downloads the provider asset into the business workspace and records provenance,
byte count and SHA-256 evidence. A manual status button remains available, but
is no longer required for an unattended approved render to complete.

Physical verification produced a two-second MP4 with the following properties:

\begin{itemize}
    \item $1080\times1920$ pixels and $9{:}16$ display aspect ratio;
    \item H.264 High Profile video;
    \item AAC stereo audio at 48~kHz;
    \item 30 frames per second;
    \item locally stored thumbnail; and
    \item SHA-256 digest
    \texttt{80b1265eba012e690adb15cd833976f46e6dbca2d52e79dd0c7f2f5409a24b54}.
\end{itemize}

The verified output is stored at
\path{/Users/kevinrobinson/Tessaris/workspaces/demo/marketing/timelines/timeline_050a3f963e1545f2/instagram_reels_export.mp4}.

A second installed-application verification produced a 5.5-second two-clip
vertical export containing a fade transition, two burned subtitle cues, a timed
campaign overlay, a logo and mixed music. The service also proved that an
eight-second default trim is clamped safely to each source file's real duration.
The final visual-QA asset has SHA-256 digest
\texttt{9e8d813556c4af1198bbd416a5ef4a09bd0d51d38e7ef3de39ac54b069fb5804}.

\subsection{Premium motion, narration and publication authority}

Each clip now carries an optional static, slow zoom-in, slow zoom-out, left-pan
or right-pan motion treatment. Motion is calculated against the real source
duration and rendered locally. Titles support no animation, fade or slide-up,
and the user may select clean editorial, bold commercial or cinematic brand
treatment.

The audio graph accepts an independent voice-over track with start-time and gain
controls. When enabled, side-chain compression lowers the background music under
narration. Music receives bounded fades and the complete mix can be normalised to
the social or broadcast loudness preset. No cloud service is required for these
operations.

The desktop can also record narration directly from the authorised microphone.
The browser recorder sends the resulting Opus/WebM asset only to the local
Tessaris backend, which stores it under the current business workspace with byte
count and SHA-256 evidence before it enters the same mastering graph.

A governed publication package freezes the exact finished-asset digest, caption,
call to action, channel, organic or paid mode, schedule and budget. Exact approval
is recorded separately. The package reports missing Meta, Google Ads or platform
credentials and has no external execution endpoint in this increment, so neither
publication nor advertising spend can occur accidentally.

Campaign evidence records impressions, clicks, qualified actions, conversions,
spend, revenue, complaints and false or unsafe claims. Tessaris derives click
rate, qualified-action rate, conversion rate, cost per qualified action and
return on advertising spend. Unsafe evidence changes the package state to an
escalated measurement rather than a winner.

The premium physical proof produced a 5.4-second $1080\times1920$ H.264/AAC
master containing two motion-treated clips, dissolve, animated title, subtitles,
logo, music, separate narration, automatic ducking and social loudness
mastering. Its SHA-256 digest is
\texttt{816264199101da6915bb50dd192b963bc174e2ba755ff30766171aa4f85d5e4b}.

\subsection{HTTP interfaces}

The provider-neutral planning and experimentation boundary exposes:

\begin{itemize}
    \item \path{GET /api/aion/marketing/creative/capabilities}
    \item \path{GET /api/aion/marketing/creative/platforms}
    \item \path{POST /api/aion/marketing/creative/plan}
    \item \path{POST /api/aion/marketing/creative/experiment/evaluate}
\end{itemize}

The paid-render and local-timeline boundary exposes:

\begin{itemize}
    \item \path{GET /api/aion/marketing/creative/{workspace_id}/renders}
    \item \path{POST /api/aion/marketing/creative/{workspace_id}/renders}
    \item \path{POST /api/aion/marketing/creative/{workspace_id}/renders/{job_id}/approve}
    \item \path{POST /api/aion/marketing/creative/{workspace_id}/renders/{job_id}/submit}
    \item \path{POST /api/aion/marketing/creative/{workspace_id}/renders/{job_id}/poll}
    \item \path{POST /api/aion/marketing/creative/{workspace_id}/timelines}
    \item \path{POST /api/aion/marketing/creative/{workspace_id}/timelines/{timeline_id}/export}
\end{itemize}

The routers are mounted in both the historical full backend and the lean desktop
backend. The latter is the runtime used by the installed macOS application.

\subsection{Principal implementation files}

\begin{itemize}
    \item Desktop interface and department routing:
    \path{desktop/mac/src/app.js}
    \item Native video-picker boundary:
    \path{desktop/mac/electron/main.js} and
    \path{desktop/mac/electron/preload.js}
    \item Creative HTTP API:
    \path{backend/modules/aion_business/api/marketing_creative_api.py}
    \item Campaign and experiment engine:
    \path{backend/modules/aion_business/runtime/marketing_creative_engine.py}
    \item Provider request contracts:
    \path{backend/modules/aion_business/runtime/creative_provider_adapters.py}
    \item Persistent paid-render service:
    \path{backend/modules/aion_business/runtime/creative_render_service.py}
    \item Automatic provider completion monitor:
    \path{backend/modules/aion_business/runtime/creative_render_polling_service.py}
    \item Local timeline and export service:
    \path{backend/modules/aion_business/runtime/creative_timeline_service.py}
    \item Publication authority and campaign evidence service:
    \path{backend/modules/aion_business/runtime/creative_publication_service.py}
    \item Lean desktop application boundary:
    \path{backend/desktop_app.py}
\end{itemize}

\subsection{Verification and recovery}

The department-shell standardisation passed 37 relevant regression checks. The
Marketing consolidation passed 65 focused checks. Additional focused checks
cover creative-plan construction, provider request contracts, evidence thresholds,
truthfulness safeguards, paid-render approval and timeline export.

The installed desktop backend reports healthy in lean
\texttt{tessaris\_desktop} mode. Relevant recoverable application snapshots are:

\begin{itemize}
    \item \path{/Applications/Tessaris.app.backup-20260810-before-support-layout-consolidation}
    \item \path{/Applications/Tessaris.app.backup-20260810-before-unified-department-shells}
    \item \path{/Applications/Tessaris.app.backup-20260810-before-unified-marketing-studio}
    \item \path{/Applications/Tessaris.app.backup-20260810-before-creative-intelligence-lab}
    \item \path{/Applications/Tessaris.app.backup-20260810-before-paid-render-and-local-timeline}
    \item \path{/Applications/Tessaris.app.backup-20260810-before-department-pilots-sidebar-link}
    \item \path{/Applications/Tessaris.app.backup-20260810-before-marketing-editor-completion}
    \item \path{/Applications/Tessaris.app.backup-20260810-before-premium-marketing-production}
    \item \path{/Applications/Tessaris.app.backup-20260810-before-local-narration-recorder}
\end{itemize}

No Git commit was made during these increments.

\subsection{Honest remaining boundary}

The architecture, governed render execution and physical platform export are now
operational. The remaining work is primarily production-channel completion and a
richer editor rather than another foundational redesign:

\begin{itemize}
    \item free-form keyframe curves beyond the safe motion presets and fine-grained caption typography;
    \item provider-backed speech generation as an optional alternative to imported or locally recorded narration;
    \item provider-native webhooks as an optimisation over the operational automatic poller;
    \item optional Runway, Luma and fal credential/execution adapters;
    \item executable Meta, Google and other publishing or advertising connectors after business credentials are supplied;
    \item automatic conversion feedback from Sales and imported platform analytics; and
    \item real campaign validation across multiple businesses and channels.
\end{itemize}

The intended operating rule remains: Tessaris may plan, generate, edit and learn
broadly, but paid generation, publishing, advertising spend and winner promotion
remain bounded by explicit business authority and recorded evidence.


\section{Unified Department Pilots and Marketing Creative Production Stack}
\label{sec:tessaris-department-pilots-marketing-creative}

\subsection{Scope and design objective}

This development phase consolidated the Tessaris department workspaces and
advanced Marketing from an early social-post prototype into a governed creative
production and experimentation system. The work covered five connected layers:

\begin{enumerate}
    \item a consistent department page structure for Marketing, Sales, Finance,
    Operations, Support and Human Resources;
    \item department-scoped AION/Pilot conversations and Boardroom assignment
    packages;
    \item consolidation of the former standalone Marketing Stream into the
    Marketing Pilot workspace;
    \item a provider-neutral Creative Intelligence Lab for images, video, clips,
    platform renditions and controlled experiments; and
    \item exact-approved paid video rendering plus locally owned timeline and
    clip export.
\end{enumerate}

The central architectural principle is that Tessaris owns the business context,
creative direction, governance, asset provenance and learning loop. External
generative-AI providers remain replaceable production engines. This allows the
product to perform as much work as possible locally while using specialist or
premium services only where model-scale generation is genuinely required.

\subsection{Unified department workspace shell}

The earlier department pages had diverged. Some routes rendered a full-screen
spatial Boardroom, some retained an obsolete black Pilot terminal, and some hid
the newer operational workspace. Support and Human Resources were especially
affected because legacy containers could overlap or replace the newer interfaces.

Every principal department now follows the same hierarchy:

\begin{enumerate}
    \item \textbf{Compact department director.} A small department-specific AION
    character identifies the function and its advisory boundary. It replaces the
    oversized spatial room on operational pages.
    \item \textbf{Department-scoped conversation.} The \emph{Talk to Pilot}
    control opens a conversation whose context is limited to the active
    department. For example, the Support Pilot discusses customer cases,
    policies and escalations rather than acting as a generic assistant.
    \item \textbf{Boardroom assignment package.} The package is retained because
    it contains the Boardroom meeting instructions, goals, discovery requirements,
    linked goal sheet and expected outputs. The \emph{Open Package} control is
    available consistently.
    \item \textbf{Operational workspace.} The real department system appears
    beneath the package: Sales Revenue, Finance, Support Case Centre, HR people
    and permissions, Marketing, or the Operations foundation.
\end{enumerate}

Legacy black Pilot containers and oversized spatial canvases are explicitly
excluded from these routes when they would override the operational workspace.
The installed department outcomes are:

\begin{table}[htbp]
\centering
\small
\begin{tabular}{p{0.17\linewidth} p{0.70\linewidth}}
\toprule
\textbf{Department} & \textbf{Mounted operational workspace} \\
\midrule
Marketing & Unified automation, creative production, approvals, calendar and history \\
Sales & Canonical CRM, governed communications and Sales Revenue workspace \\
Finance & Finance operating workspace, evidence, bookkeeping and reporting controls \\
Operations & Unified department shell and foundation for the Operations build \\
Support & Canonical Support Case Centre, policy authority and escalation workflow \\
Human Resources & Organisation, people, roles, permissions and authority workspace \\
\bottomrule
\end{tabular}
\caption{Standardised department workspace structure.}
\end{table}

The sidebar now contains a clearly labelled \textbf{Department Pilots} route,
represented by the letter \textbf{P} when the rail is collapsed. It returns the
user to the department workspace containing Marketing, Sales, Finance,
Operations, Support and Human Resources.

\subsection{Marketing workspace consolidation}

Marketing previously existed in two disconnected locations: the Marketing Pilot
department and a separate Marketing Stream in the left sidebar. The Stream
contained the early manual content form, run history, approvals and calendar,
while the Pilot could produce an isolated artefact without launching the real
Marketing workflow.

Marketing now has one canonical home with the following components:

\begin{description}
    \item[AION Automation Studio:] The user describes a business outcome. AION
    reads the Boardroom assignment, Brand Foundation, tone of voice, target
    customer, offers and channel rules before launching the governed Marketing
    workflow.
    \item[Manual Creation Studio:] The user can directly specify the brief,
    objective, funnel goal, audience, persona, offer, channels, keywords,
    hashtags, hard rules, reference images, product assets and creative direction.
    \item[Shared workflow:] Automated and manual work use the same Brand
    Foundation, approval queue, content-generation path, run history and output
    store.
    \item[Operational navigation:] Stream, Calendar, Runs, Approvals and Settings
    remain accessible inside Marketing.
\end{description}

The obsolete standalone Marketing Stream link was removed. Historical routes are
redirected into Marketing, and \emph{Build with AION} now launches the real
Marketing workflow rather than producing an unconnected front-end artefact.

The consolidated foundation retains caption drafting, five-card carousel copy,
product and reference-image intake, the existing OpenAI image route, social-post
packaging, approvals, run history and failure evidence.

\subsection{Creative Intelligence Lab}

The Creative Intelligence Lab adds campaign-level creative orchestration above
the existing Marketing workflow. It converts one approved business idea into a
master concept, storyboard, generation jobs, platform renditions and a measurable
experiment.

\subsubsection{Creative planning}

For each campaign, the system records or derives:

\begin{itemize}
    \item campaign promise, objective, target audience and approved offer;
    \item tone, proof requirements and call to action;
    \item a five-scene structure covering hook, context, solution, proof and call
    to action;
    \item image, original-video, short-clip and carousel requirements;
    \item shot direction, captions, subtitle contracts and visual safe zones; and
    \item a stable plan identifier and SHA-256-based plan hash.
\end{itemize}

Generated claims must be supported by recorded business evidence. The engine is
not permitted to invent reviews, outcomes, prices or product capabilities.

\subsubsection{Platform renditions}

The platform-profile registry produces channel-specific rather than merely
resized outputs. Profiles currently cover Instagram feed, Reels and Stories;
Facebook feed and Reels; TikTok; YouTube Shorts; and LinkedIn feed and video.
The principal production contracts are:

\begin{itemize}
    \item $4{:}5$ portrait feed assets at $1080\times1350$ pixels;
    \item $9{:}16$ vertical video and clip assets at $1080\times1920$ pixels;
    \item platform-appropriate duration guidance;
    \item essential text, faces, logos and calls to action inside the central safe
    area; and
    \item burned-in caption guidance plus an editable subtitle-track contract for
    spoken video.
\end{itemize}

\subsubsection{Provider strategy}

Tessaris separates the creative operating system from the rendering model:

\begin{table}[htbp]
\centering
\small
\begin{tabular}{p{0.20\linewidth} p{0.28\linewidth} p{0.39\linewidth}}
\toprule
\textbf{Layer} & \textbf{Providers} & \textbf{Responsibility} \\
\midrule
Primary direct engines & OpenAI GPT Image/Sora; Google Imagen/Veo & Image and source-video generation through existing provider credentials \\
Optional specialists & Runway, Luma and fal & Premium or specialised models when directly connected and approved \\
Product references & Higgsfield and OpenArt & User-experience and capability references; not mandatory dependencies \\
Tessaris native layer & Tessaris/AION & Brand context, storyboards, prompts, platform adaptation, clipping, approvals, evidence and learning \\
\bottomrule
\end{tabular}
\caption{Provider-neutral creative-production architecture.}
\end{table}

This design prevents core Marketing capability from becoming locked to a
third-party wrapper whose models, prices or API availability may change.

\subsection{Controlled experimentation and winner learning}

The experiment governor creates a control and up to three challengers. Each
challenger changes exactly one variable:

\begin{itemize}
    \item the opening hook;
    \item the opening visual; or
    \item the call to action.
\end{itemize}

The default evidence contract requires at least 100 observations per variant, a
minimum relative improvement of 15\%, and no breach of complaint or truthfulness
guardrails. Challenger traffic is capped at 20\% until a result is accepted.

The evaluator can return four meaningful outcomes:

\begin{description}
    \item[Collect more evidence:] One or more variants have not reached the frozen
    sample threshold.
    \item[Keep control:] No challenger clears the improvement threshold, or an
    apparent improvement worsens a protected metric such as complaints.
    \item[Stop and review:] A false-claim, unsafe-content or equivalent guardrail
    is breached.
    \item[Winner candidate:] A challenger clears the primary threshold without
    breaching a guardrail. It still requires approval before traffic changes.
\end{description}

No evaluation automatically publishes content, purchases advertising, changes
campaign spend or promotes a winner. A live verification used simulated metrics
and correctly identified a challenger with a 30\% improvement as a candidate
while keeping automatic application disabled.

\subsection{Paid source-video rendering}

The paid-render layer implements persistent asynchronous jobs for OpenAI Sora and
Google Veo. A render job stores:

\begin{itemize}
    \item business workspace and campaign-plan reference;
    \item provider and exact model;
    \item full prompt and optional reference-asset path;
    \item aspect ratio, duration and resolution;
    \item rendition identifier and immutable approval hash;
    \item provider job identifier and provider response;
    \item completion status, error state and timestamps; and
    \item locally owned output metadata.
\end{itemize}

The lifecycle is deliberately separated into the following operations:

\begin{enumerate}
    \item \textbf{Prepare.} Build and persist the complete render contract without
    making a paid provider request.
    \item \textbf{Approve.} Bind the unchanged contract to an identified approver,
    an explicit paid-generation boolean and a positive maximum euro cost.
    \item \textbf{Submit.} Recalculate and compare the approval hash before calling
    the selected provider.
    \item \textbf{Poll.} Read the asynchronous provider status and fail closed on
    provider errors.
    \item \textbf{Acquire.} Download completed footage into the relevant Tessaris
    business workspace rather than relying on an expiring provider URL.
\end{enumerate}

A paid request without the explicit authority boolean returns HTTP status 403.
Verification prepared a dummy render contract and approved a maximum cost of
EUR~5, but deliberately did not submit the provider job; therefore no generation
charge was incurred.

Completed source footage is recorded with its local path, MIME type, byte count,
SHA-256 digest and an explicit \texttt{owned\_locally} flag. Provider credentials
are read at execution time and are never stored in render-job records.

\subsection{Native timeline, clipping and export}

Basic post-production is performed locally using the FFmpeg 7.1 binary already
available in the Tessaris Python environment. This avoids paying a premium
provider for deterministic editing tasks.

A timeline project contains an ordered set of source clips with start and end
times, labels and a target platform. Export performs the following operations:

\begin{itemize}
    \item validates every source path and trim boundary;
    \item trims and normalises each segment;
    \item scales while preserving aspect ratio;
    \item pads to the exact platform frame and produces square pixels;
    \item converts video to H.264 at 30 frames per second;
    \item preserves source audio as AAC stereo at 48~kHz;
    \item generates a silent audio stream when a source contains no audio;
    \item concatenates ordered segments and enables MP4 fast-start;
    \item creates a JPEG thumbnail; and
    \item stores output dimensions, byte count and SHA-256 evidence.
\end{itemize}

The desktop interface uses a minimal Electron inter-process bridge to open the
native macOS video picker. The renderer receives only the selected file path;
arbitrary filesystem browsing is not exposed to webpage code.

The completed visual editor extends this boundary with multi-select video,
image and audio pickers. Clips appear as draggable cards and retain independent
in/out points. The editor records bounded cuts, fades, dissolves, wipes and
slides; parses editable subtitle cues written as
\texttt{start-end|text}; applies timed text overlays; composites a local logo;
and loops and mixes background music at a user-selected volume. FFmpeg performs
all deterministic composition locally. Uploaded footage is not sent to a third
party by the editing operation.

Reference-image handling is provider aware. OpenAI jobs accept at most one
reference image, while Gemini jobs accept at most three. The selected paths form
part of the immutable approval hash. Gemini references are encoded into the
documented inline-data structure and reference-guided jobs are fixed to the
supported eight-second duration. No paid provider request occurs merely because
an image is selected.

Submitted render jobs are monitored by a daemon background poller. Completion
downloads the provider asset into the business workspace and records provenance,
byte count and SHA-256 evidence. A manual status button remains available, but
is no longer required for an unattended approved render to complete.

Physical verification produced a two-second MP4 with the following properties:

\begin{itemize}
    \item $1080\times1920$ pixels and $9{:}16$ display aspect ratio;
    \item H.264 High Profile video;
    \item AAC stereo audio at 48~kHz;
    \item 30 frames per second;
    \item locally stored thumbnail; and
    \item SHA-256 digest
    \texttt{80b1265eba012e690adb15cd833976f46e6dbca2d52e79dd0c7f2f5409a24b54}.
\end{itemize}

The verified output is stored at
\path{/Users/kevinrobinson/Tessaris/workspaces/demo/marketing/timelines/timeline_050a3f963e1545f2/instagram_reels_export.mp4}.

A second installed-application verification produced a 5.5-second two-clip
vertical export containing a fade transition, two burned subtitle cues, a timed
campaign overlay, a logo and mixed music. The service also proved that an
eight-second default trim is clamped safely to each source file's real duration.
The final visual-QA asset has SHA-256 digest
\texttt{9e8d813556c4af1198bbd416a5ef4a09bd0d51d38e7ef3de39ac54b069fb5804}.

\subsection{Premium motion, narration and publication authority}

Each clip now carries an optional static, slow zoom-in, slow zoom-out, left-pan
or right-pan motion treatment. Motion is calculated against the real source
duration and rendered locally. Titles support no animation, fade or slide-up,
and the user may select clean editorial, bold commercial or cinematic brand
treatment.

The audio graph accepts an independent voice-over track with start-time and gain
controls. When enabled, side-chain compression lowers the background music under
narration. Music receives bounded fades and the complete mix can be normalised to
the social or broadcast loudness preset. No cloud service is required for these
operations.

The desktop can also record narration directly from the authorised microphone.
The browser recorder sends the resulting Opus/WebM asset only to the local
Tessaris backend, which stores it under the current business workspace with byte
count and SHA-256 evidence before it enters the same mastering graph.

A governed publication package freezes the exact finished-asset digest, caption,
call to action, channel, organic or paid mode, schedule and budget. Exact approval
is recorded separately. The package reports missing Meta, Google Ads or platform
credentials and has no external execution endpoint in this increment, so neither
publication nor advertising spend can occur accidentally.

Campaign evidence records impressions, clicks, qualified actions, conversions,
spend, revenue, complaints and false or unsafe claims. Tessaris derives click
rate, qualified-action rate, conversion rate, cost per qualified action and
return on advertising spend. Unsafe evidence changes the package state to an
escalated measurement rather than a winner.

The premium physical proof produced a 5.4-second $1080\times1920$ H.264/AAC
master containing two motion-treated clips, dissolve, animated title, subtitles,
logo, music, separate narration, automatic ducking and social loudness
mastering. Its SHA-256 digest is
\texttt{816264199101da6915bb50dd192b963bc174e2ba755ff30766171aa4f85d5e4b}.

\subsection{HTTP interfaces}

The provider-neutral planning and experimentation boundary exposes:

\begin{itemize}
    \item \path{GET /api/aion/marketing/creative/capabilities}
    \item \path{GET /api/aion/marketing/creative/platforms}
    \item \path{POST /api/aion/marketing/creative/plan}
    \item \path{POST /api/aion/marketing/creative/experiment/evaluate}
\end{itemize}

The paid-render and local-timeline boundary exposes:

\begin{itemize}
    \item \path{GET /api/aion/marketing/creative/{workspace_id}/renders}
    \item \path{POST /api/aion/marketing/creative/{workspace_id}/renders}
    \item \path{POST /api/aion/marketing/creative/{workspace_id}/renders/{job_id}/approve}
    \item \path{POST /api/aion/marketing/creative/{workspace_id}/renders/{job_id}/submit}
    \item \path{POST /api/aion/marketing/creative/{workspace_id}/renders/{job_id}/poll}
    \item \path{POST /api/aion/marketing/creative/{workspace_id}/timelines}
    \item \path{POST /api/aion/marketing/creative/{workspace_id}/timelines/{timeline_id}/export}
\end{itemize}

The routers are mounted in both the historical full backend and the lean desktop
backend. The latter is the runtime used by the installed macOS application.

\subsection{Principal implementation files}

\begin{itemize}
    \item Desktop interface and department routing:
    \path{desktop/mac/src/app.js}
    \item Native video-picker boundary:
    \path{desktop/mac/electron/main.js} and
    \path{desktop/mac/electron/preload.js}
    \item Creative HTTP API:
    \path{backend/modules/aion_business/api/marketing_creative_api.py}
    \item Campaign and experiment engine:
    \path{backend/modules/aion_business/runtime/marketing_creative_engine.py}
    \item Provider request contracts:
    \path{backend/modules/aion_business/runtime/creative_provider_adapters.py}
    \item Persistent paid-render service:
    \path{backend/modules/aion_business/runtime/creative_render_service.py}
    \item Automatic provider completion monitor:
    \path{backend/modules/aion_business/runtime/creative_render_polling_service.py}
    \item Local timeline and export service:
    \path{backend/modules/aion_business/runtime/creative_timeline_service.py}
    \item Publication authority and campaign evidence service:
    \path{backend/modules/aion_business/runtime/creative_publication_service.py}
    \item Lean desktop application boundary:
    \path{backend/desktop_app.py}
\end{itemize}

\subsection{Verification and recovery}

The department-shell standardisation passed 37 relevant regression checks. The
Marketing consolidation passed 65 focused checks. Additional focused checks
cover creative-plan construction, provider request contracts, evidence thresholds,
truthfulness safeguards, paid-render approval and timeline export.

The installed desktop backend reports healthy in lean
\texttt{tessaris\_desktop} mode. Relevant recoverable application snapshots are:

\begin{itemize}
    \item \path{/Applications/Tessaris.app.backup-20260810-before-support-layout-consolidation}
    \item \path{/Applications/Tessaris.app.backup-20260810-before-unified-department-shells}
    \item \path{/Applications/Tessaris.app.backup-20260810-before-unified-marketing-studio}
    \item \path{/Applications/Tessaris.app.backup-20260810-before-creative-intelligence-lab}
    \item \path{/Applications/Tessaris.app.backup-20260810-before-paid-render-and-local-timeline}
    \item \path{/Applications/Tessaris.app.backup-20260810-before-department-pilots-sidebar-link}
    \item \path{/Applications/Tessaris.app.backup-20260810-before-marketing-editor-completion}
    \item \path{/Applications/Tessaris.app.backup-20260810-before-premium-marketing-production}
    \item \path{/Applications/Tessaris.app.backup-20260810-before-local-narration-recorder}
\end{itemize}

No Git commit was made during these increments.

\subsection{Honest remaining boundary}

The architecture, governed render execution and physical platform export are now
operational. The remaining work is primarily production-channel completion and a
richer editor rather than another foundational redesign:

\begin{itemize}
    \item free-form keyframe curves beyond the safe motion presets and fine-grained caption typography;
    \item provider-backed speech generation as an optional alternative to imported or locally recorded narration;
    \item provider-native webhooks as an optimisation over the operational automatic poller;
    \item optional Runway, Luma and fal credential/execution adapters;
    \item executable Meta, Google and other publishing or advertising connectors after business credentials are supplied;
    \item automatic conversion feedback from Sales and imported platform analytics; and
    \item real campaign validation across multiple businesses and channels.
\end{itemize}

The intended operating rule remains: Tessaris may plan, generate, edit and learn
broadly, but paid generation, publishing, advertising spend and winner promotion
remain bounded by explicit business authority and recorded evidence.
